\fancysection{Expected Outcomes}
%
\label{sec:expected_outcomes}

\subsection{Catalog Size}
%
\label{sec:catalog_size}

\citep{Prsa_et_al_22} searched $\sim2\times10^5$ LCs extracted from the high
cadence stamp of the first 26 \tess{} sectors, detecting 4\,584 EBs. Not all of
these binaries will be included in the catalog we plan to produce since the
\citep{Prsa_et_al_22} catalog includes contact binaries and binaries for which
the signal to noise will not be sufficient for high precision physical parameter
determination.  Requiring detached binaries (the sum of the stellar radii not
exceeding half the pericenter distance), primary eclipse depth of at least 5\%,
secondary eclipse depth of at least 1\% allows us to estimate that we will be
able to obtain high quality and high precision characterization of at least 800
binaries from stamp LCs with about double that number reasonably well
constrained. We obtain the same estimate by scaling the results of
\citep{Justesen_Albrecht_21} who found ~350 high precision binaries and more
than 800 with reasonable constraints just from the southern hemisphere. With
double or triple the observations for each binary, we will undoubtedly be able
to extract high precision physical parameters from somewhat shallower eclipses
than \citep{Justesen_Albrecht_21}, and will have improved sensitivity to longer
orbital periods. Therefore, \textbf{to a good approximation our EB catalog based
on stamp data alone should include $\sim1500$ EBs, roughly half of which will
have high precision measurements of their parameters.}

Predicting the expected EB detections from FFIs is less straightforward. We plan
to use the Quick Look Pipeline (QLP) LCs \citep{Huang_et_al_20a} based on the
full frames. The QLP provides LCs for more than $10^7$ stars brighter than
\tess{} magnitude $T<13.5$. Naively, we will search roughly 50 times more LCs
and the 30 min cadence is still plenty to resolve all features of the LCs. As a
result, one could estimate creating a catalog of $5\times10^4$ EBs, half with
high precision parameters. However, the $2\times10^5$ stamps targets were not
randomly selected. In fact over 1700 known or suspected EBs were specifically
added to the target list through the TESS Guest Investigator program. On the
other hand, the QLP LC database only includes stars down to $T<13.5$ whereas
stamp targets included much fainter stars. The photometric precision on a 1-h
timescale vs magnitude reported for the stamp-based and FFI-based LCs is almost
identical: compare for example \citep{Kunimoto_et_al_22} to the \tess{}
instrument handbook \citep{TESS_handbook}. Hence, on average the FFI LCs will
have higher precision than stamp LCs. Regardless, it is reasonable to expect
that \textbf{we will end up with $\sim10^4$ EBs with high precision parameters
and roughly double to several times that number in total from the FFI LCs.}

In order to get a lower limit to the fraction of binaries for which we also
expect to measure the spin period and often differential rotation, we look to
\citep{CantoMartins_et_al_20}, who searched 1000 \tess{} objects of interest and
measured 131 unambiguous rotation periods (13\%) using less than a year of
observations, and \citep{Oelkers_et_al_18} who searched $\approx4\times10^6$
ground based LCs (from the KELT survey) detecting 47\,174 (12\%) with clear spin
signatures with periods under 13 days (a conservative limit to what can be
detected using \tess{} for targets where consecutive sectors do not overlap).
The former survey relied on 2 or 3 times fewer observations per target than we
will use and the latter used ground based data of much lower photometric
precision, so we can confidently conclude 10\%-15\% is a lower limit to the
fraction of stars exhibiting rotational modulations detectable by \tess{} at any
given time. When focusing on EBs that fraction is bound to be higher, since
tidal interactions spin-up short period binaries increasing the fraction with
periods under 13 days, and likely increasing the signal since faster rotators
tend to have higher magnetic activity. For an upper limit to the fraction of EBs
for which we will measure the spin, we turn to \citep{Lurie_et_al_17} who used
\kepler{} photometry (higher precision and longer time span than we will have)
to measure spin of EBs.  Crossmatching the \citep{Lurie_et_al_17} catalog of
firmly detected spins with the \citep{Windemuth_et_al_19} catalog of EBs with
high precision parameters results in 32\% (233 out of 728 EBs) of the
\citep{Windemuth_et_al_19} stars having a spin period less than 13 days. In
short, \textbf{we will be able to detect spin periods for between 10\% and 30\%
of the EBs in our catalog: several hundred from stamps and many thousands from
FFIs.}

\subsection{Expected Precision}
%
\label{sec:expected_precision}

In addition to the number of systems, it would be useful to have an idea of the
precision with which various physical parameters will be determined. Due to the
generally highly correlated uncertainties of the parameters, the precision that
is relevant to one application may be very different from another. For example,
tides are sensitive to the ratio of stellar radii to the semimajor axis. This is
usually much better constrained than either of the radii or the semimajor axis
by itself (marginalized over all other quantities). Regardless, marginalized
uncertainties of individual parameters can still be a useful guide to the value
of the catalog. While \citep{Justesen_Albrecht_21} did not publish
uncertainties, since they were not able to run their analysis long enough to be
fully converged, the order of magnitude of the marginalized uncertainties can be
estimated from the available samples. \textbf{For the high precision EBs
mentioned above, we expect marginalized uncertainties for the stellar masses
better than 3\%, and for the radii better than 2\% \citep[private
communication]{Justesen_Albrecht_21}. For eccentricities, the $e\cos\omega$
component is much better constrained, with typical uncertainties
$\mathbf{\sim10^{-4}}$, than the $e\sin\omega$ component, with typical
uncertainties $\mathbf{\sim10^{-3}}$ \citep{Justesen_Albrecht_21}.} Measuring
spin periods from spot modulation usually produces extremely precise values.
However, the real uncertainty in that quantity comes from the fact that
different longitudes of the same star rotate at different rates. As a result,
the uncertainty in the location of the stellar spots producing the rotational
modulation dominates the uncertainty with which the rotation period can be
measured. As a result, \textbf{the spin precision will be the best theoretically
possible}, but still limited to $\sim$10\% due to differential rotation.

The precision with which eccentricity is determined is matched only by the much
smaller catalog of EBs from \kepler{} \citep{Windemuth_et_al_19}. That level of
precision is crucial in tidal studies as it allows us to probe the final stages
of tidal circularization, where tidal deformations are simplest (a single tidal
frequency dominates by a very large margin over all others), which in turn makes
for much easier and more reliable interpretation of the inferred tidal
dissipation.
